\section{Conclusion}
\label{sec:conclusion}

Structured quantum tensor networks invite a simulation analysis that uses both graph structure and stabilizer resource theory. Treewidth bounds capture separators but ignore whether the tensors crossing them are free or resourceful. Global magic bounds capture non-stabilizer content but ignore where that content sits in a compositional graph.

The contribution is a certificate-based separator-local theorem for locally marginalizable resource-decomposition networks. The key parameter is \(\Lammsg\), computed by an explicit active-child recurrence, together with the stabilizer-compressed width \(\weff\). The scalar sampling dependence is \(\exp(2\Lammsg)\). This theorem proves a restricted easy regime and gives pedagogical separations from naive global-magic accounting on constant-boundary families with many locally consumed resources, including a path-active variant where unresolved resource randomness persists through \(O(\log n)\) separators rather than vanishing immediately at leaves. It also gives a stabilizer-grid parameter separation between dense treewidth bounds and stabilizer-compressed separator-local certificates.

The accompanying artifact makes the certificate concrete. It attempts real lambeq extraction when available, falls back to explicitly labelled structured instances otherwise, exports \(\xi_b\), \(\kappa_b\), \(\Act(b)\), \(\Lammsg\), \(\Lambda_{\mathrm{glob}}\), width proxies, and comparator exponent proxies, and records the run manifest. In the executed run reported here, lambeq imported successfully in a Python 3.11 environment, but no Bobcat/DepCCG/CCG parser produced parsed diagrams for the sample corpus. A local lambeq LinearReader path did produce 120 graph-only diagrams with no resource annotations, and the run therefore used 120 fallback structured instances as the primary nontrivial certificate evidence, with 105 certificate passes and 15 deliberate certificate failures. The artifact also includes small dense-contraction and scalar-estimator microbenchmarks on graph-derived data. The generated plots and CSV files are certificate-evidence diagnostics rather than empirical parser-level QNLP results.

SLMS should be viewed as a sufficient-condition framework for identifying classically easy structured quantum tensor networks. Its value is diagnostic: it makes explicit when local non-stabilizer resources are consumed before loading high-level separators, and records this in a reproducible certificate. The present artifact demonstrates the mechanics of the certificate, its failure modes, and its comparison to several implemented proxy diagnostics. Extending the same workflow to parser-generated QNLP/QDisCoCirc resource instances and to exact ZX/STN implementations remains future work.
